%% kapa: exam=1 round=36 year=2025 period=1 form=A subject=civil_law subject_ko=민법 seq=11 number=11
%% source: Q-Net 감정평가사 기출문제 (cst00309) artlSeq=5231525 file=raw/1008/5231525/2025년 제36회 감정평가사 1차시험 1교시 문제지 원본.pdf
%% conversion: kordoc 4.24.51 -> tools/gen1.py (auto); review=proofread-sonnet; flags=none
\begin{question}{11}
진의 아닌 의사표시에 관한 설명으로 옳지 않은 것은? (다툼이 있으면 판례에 따름)
\begin{choices}
  \item 진의 아닌 의사표시에 관한 민법 제107조는 공무원의 사직의 의사표시에는 적용되지 않는다.
  \item 진의 아닌 의사표시에서 '진의'란 특정한 내용의 의사표시를 하고자 하는 표의자의 생각을 말한다.
  \item 진의 아닌 의사표시의 무효를 주장하는 경우, 상대방의 악의 또는 과실은 그 주장자가 증명해야 한다.
  \item 대리행위에서 진의 아닌 의사표시인지 여부는 대리인을 표준으로 정한다.
  \item 표의자가 강박에 의하여 증여를 하기로 하고 그에 따른 증여의 의사표시를 하였더라도 재산을 강제로 뺏긴다는 것이 표의자의 본심으로 잠재되어 있었다면 증여의 내심의 효과의사가 결여된 것이다.
\end{choices}
\end{question}
\answer{5}
